\setcounter{table}{0}
\begin{table}[!ht]\centering
\caption{Studies of body-part or anatomy classification, and what each reports}\label{tab:literature}
{\footnotesize\setlength{\tabcolsep}{3pt}\begin{tabular}{@{}>{\raggedright\arraybackslash}p{0.15\textwidth}>{\raggedright\arraybackslash}p{0.19\textwidth}>{\raggedright\arraybackslash}p{0.10\textwidth}>{\raggedright\arraybackslash}p{0.09\textwidth}>{\raggedright\arraybackslash}p{0.13\textwidth}>{\raggedright\arraybackslash}p{0.12\textwidth}>{\raggedright\arraybackslash}p{0.08\textwidth}@{}}
\toprule
\textbf{study} & \textbf{images and classes} & \textbf{data} & \textbf{split unit} & \textbf{label review} & \textbf{external data} & \textbf{calib.} \\
\midrule
Pham et al. (2021) & 16,093 radiographs, 5 coarse groups & public & - & - & - & - \\
Langer et al. (2024, preprint) & 48,434 radiographs, 14 regions & hospital & - & experts reviewed test errors & - & - \\
Mitsuyama et al. (2026) & 429,341 radiographs, 7 categories & 2 sites & - & radiologists corrected tags & second institution, MURA & - \\
Krishnan and Padmavathi (2025) & UNIFESP, 22 classes, single label & public & not described & - & - & - \\
Shah et al. (2026) & MURA, 7 study types & public & patient-wise & - & - & - \\
Fink et al. (2024) & musculoskeletal radiographs, projection and side & - & patient & - & - & - \\
Raffy et al. (2023) & CT and MRI, 17 and 18 regions & - & - & - & 27 institutions & - \\
This work & UNIFESP 22 classes (multilabel), MURA 7 classes, FracAtlas, COVID-19 & public only & patient, nested & audit of 39 suspects (no clinician) & MURA, FracAtlas, UNIFESP & ECE, risk-coverage \\
\bottomrule
\end{tabular}}
\par\smallskip{\footnotesize calib. = calibration or selective prediction; 2 sites = two institutions. A dash means the item was not reported in the abstract or in the parts of the paper we read; we did not check every full text, so a dash is not proof of absence.}
\end{table}
